\begin{abstract}
Wireless sensor nodes spend most of their energy on the radio and rarely get a
new battery, so lifetime is decided by how the network organizes its
transmissions. Clustering is the standard answer, and twenty-five years of
clustering protocols now span three methodological generations: probabilistic
and threshold heuristics, explicit optimization, and learned selection
functions. Those generations are usually compared by reading their published
tables side by side, which does not work, because each protocol was measured
against a different simulator, field size, node count, energy budget and
channel. We rebuild the comparison from scratch. Nine protocols spanning all
three generations, plus a direct-transmission baseline, run on a single engine
that owns energy accounting, packet sizing, fusion, the channel and the
counting of delivered readings; a protocol returns only a cluster structure, and
a static audit confirms that no implementation touches the energy or liveness
arrays. Trials are paired, so at a given index every protocol sees the same
topology, the same per-link shadowing and the same sensed-value stream, and the
resulting differences are tested with sign-flip permutation, paired bootstrap
intervals and Holm correction. Running the whole study at both endpoints of the
one parameter the source specifications leave open, transmit power, retracts one
of our own comparisons: the first-death advantage of the type-2 fuzzy protocol
and the deep Q-network over LEACH is 417 and 426 rounds under packet loss and
falls to 17 and 22 rounds, no longer significant, once loss is removed. Their
advantage in overall survival and in readings delivered survives both channels.
A 1{,}350-run sweep over three node counts and three field areas shows the same
boundary from the other side: distance to the sink, not node density, sets the
regime, and in a $50\times50$~m field clustering is worse than no clustering at
all. We report the mechanism by measurement rather than by argument, and we
disclose the modeling choices that bound the results, ordered by how far each
could move a conclusion.
\end{abstract}

\begin{IEEEkeywords}
Wireless sensor networks, clustering protocols, network lifetime, energy
efficiency, reinforcement learning, graph neural networks, type-2 fuzzy logic,
cluster head selection, reproducibility.
\end{IEEEkeywords}
